% GENERATED FILE. Edit canonical lessons, then run npm run generate:books.
% canonical-source-hash: fnv1a64:034fea59f560de1a
% canonical-lessons: FA-C522-qafiye, FA-C522-sorudan, FA-C522-dekleme-kardan, FA-C522-elham-gereftan, FA-C522-takhayyol

\chapter{Writing Verse}
\label{ch:fa-writing-verse}
\hlchaptermodality{522}

\begin{chapteropening}
\textbf{By the end of this chapter:} \emph{I can talk about writing and reciting poetry.}

Complete the last lesson of chapter 522: I can talk about writing and reciting poetry.
\end{chapteropening}

\section[\texorpdfstring{\fa{قافیه} (qâfiye)}{قافیه (qâfiye)}]{\fa{قافیه} (qâfiye) — a rhyme}

\label{lesson:FA-C522-qafiye}



\begin{quote}
Before the new one: say the Persian for a quatrain, then the Persian for a couplet.
\end{quote}

\subsection*{You'll want to know: \fa{قافیه}}
\textbf{\fa{قافیه}} (\emph{qâfiye}) — \textquotedblleft{}a rhyme\textquotedblright{}.

\begin{quote}
\fa{این} \fa{دو} \fa{کلمه} \fa{با} \fa{هم} \fa{قافیه} \fa{دارند}.
\end{quote}

\begin{quote}
\emph{in do kalame bâ ham qâfiye dârand.}
\end{quote}

\begin{quote}
These two words rhyme.
\end{quote}

\emph{qâfiye dâshtan} — to rhyme. Plural \emph{qavâfi} (Arabic).

\begin{cousinweb}
\emph{qâfiye} is a loan, usually traced to Arabic \emph{qāfiya}, 'rhyme'. EtymDB has no entry for it.
\end{cousinweb}

\begin{grammarlens}[title={What you've built}]
One more word for telling, explaining and describing at B1.
\end{grammarlens}

\subsection*{Your turn}
\begin{itemize}
\raggedright
  \item \emph{Say it:} \emph{qâfiye}
  \item \emph{Say it:} \emph{qâfiye}, once more
  \item \emph{Say it:} a sentence of your own with \emph{qâfiye}
  \item \emph{Recall:} say the Persian for a quatrain, then the Persian for a couplet, then say \emph{qâfiye} again
\end{itemize}

\begin{tcolorbox}[breakable,colback=teal!4,colframe=teal!35!black,title={Before you move on}]
What does \fa{قافیه} mean? (\textquotedblleft{}A rhyme\textquotedblright{}.) Say it once more.
\end{tcolorbox}
\section[\texorpdfstring{\fa{سرودن} (sorudan)}{سرودن (sorudan)}]{\fa{سرودن} (sorudan) — to compose (poetry)}

\label{lesson:FA-C522-sorudan}



\begin{quote}
Before the new one: say the Persian for a couplet, then the Persian for a rhyme.
\end{quote}

\subsection*{You'll want to know: \fa{سرودن}}
\textbf{\fa{سرودن}} (\emph{sorudan}) — \textquotedblleft{}to compose (poetry)\textquotedblright{}.

\begin{quote}
\fa{او} \fa{این} \fa{شعر} \fa{را} \fa{در} \fa{جوانی} \fa{سرود}.
\end{quote}

\begin{quote}
\emph{u in she'r râ dar javâni sorud.}
\end{quote}

\begin{quote}
He composed this poem in his youth.
\end{quote}

Literary; present stem \emph{sorâ-}: \emph{misorâyad}. In speech people say \emph{she'r goftan}.

\begin{cousinweb}
\emph{sorudan} is a native Persian verb, with the present stem \emph{sorâ-}, but EtymDB lists it without a source, and its older history is not settled.
\end{cousinweb}

\begin{grammarlens}[title={What you've built}]
One more word for telling, explaining and describing at B1.
\end{grammarlens}

\subsection*{Your turn}
\begin{itemize}
\raggedright
  \item \emph{Say it:} \emph{sorudan}
  \item \emph{Say it:} \emph{sorudan}, once more
  \item \emph{Say it:} a sentence of your own with \emph{sorudan}
  \item \emph{Recall:} say the Persian for a couplet, then the Persian for a rhyme, then say \emph{sorudan} again
\end{itemize}

\begin{tcolorbox}[breakable,colback=teal!4,colframe=teal!35!black,title={Before you move on}]
What does \fa{سرودن} mean? (\textquotedblleft{}To compose (poetry)\textquotedblright{}.) Say it once more.
\end{tcolorbox}
\section[\texorpdfstring{\fa{دکلمه} \fa{کردن} (dekleme kardan)}{دکلمه کردن (dekleme kardan)}]{\fa{دکلمه} \fa{کردن} (dekleme kardan) — to recite (poetry)}

\label{lesson:FA-C522-dekleme-kardan}



\begin{quote}
Before the new one: say the Persian for a rhyme, then the Persian for to compose.
\end{quote}

\subsection*{You'll want to know: \fa{دکلمه} \fa{کردن}}
\textbf{\fa{دکلمه} \fa{کردن}} (\emph{dekleme kardan}) — \textquotedblleft{}to recite (poetry)\textquotedblright{}.

\begin{quote}
\fa{شاگردها} \fa{در} \fa{جشن} \fa{شعر} \fa{دکلمه} \fa{کردند}.
\end{quote}

\begin{quote}
\emph{shâgerdhâ dar jashn she'r dekleme kardand.}
\end{quote}

\begin{quote}
The pupils recited poems at the celebration.
\end{quote}

\emph{dekleme} 'recitation' + \emph{kardan}.

\begin{cousinweb}
\emph{dekleme} is a modern loan, usually traced to French \emph{déclamer}, 'to declaim'. EtymDB has no entry for it. \emph{kardan} is inherited Persian.
\end{cousinweb}

\begin{grammarlens}[title={What you've built}]
One more word for telling, explaining and describing at B1.
\end{grammarlens}

\subsection*{Your turn}
\begin{itemize}
\raggedright
  \item \emph{Say it:} \emph{dekleme kardan}
  \item \emph{Say it:} \emph{dekleme kardan}, once more
  \item \emph{Say it:} a sentence of your own with \emph{dekleme kardan}
  \item \emph{Recall:} say the Persian for a rhyme, then the Persian for to compose, then say \emph{dekleme kardan} again
\end{itemize}

\begin{tcolorbox}[breakable,colback=teal!4,colframe=teal!35!black,title={Before you move on}]
What does \fa{دکلمه} \fa{کردن} mean? (\textquotedblleft{}To recite (poetry)\textquotedblright{}.) Say it once more.
\end{tcolorbox}
\section[\texorpdfstring{\fa{الهام} \fa{گرفتن} (elhâm gereftan)}{الهام گرفتن (elhâm gereftan)}]{\fa{الهام} \fa{گرفتن} (elhâm gereftan) — to be inspired, to draw inspiration}

\label{lesson:FA-C522-elham-gereftan}



\begin{quote}
Before the new one: say the Persian for to compose, then the Persian for to recite.
\end{quote}

\subsection*{You'll want to know: \fa{الهام} \fa{گرفتن}}
\textbf{\fa{الهام} \fa{گرفتن}} (\emph{elhâm gereftan}) — \textquotedblleft{}to be inspired, to draw inspiration\textquotedblright{}.

\begin{quote}
\fa{شاعر} \fa{از} \fa{دریا} \fa{الهام} \fa{گرفت}.
\end{quote}

\begin{quote}
\emph{shâ'er az daryâ elhâm gereft.}
\end{quote}

\begin{quote}
The poet was inspired by the sea.
\end{quote}

\emph{az chizi elhâm gereftan}; \emph{elhâm} 'inspiration' + \emph{gereftan} 'to take'.

\begin{cousinweb}
\emph{elhâm} was borrowed from Arabic \emph{ilhām}, 'inspiration' (EtymDB). \emph{gereftan}, 'to take', is inherited Persian.
\end{cousinweb}

\begin{grammarlens}[title={What you've built}]
One more word for telling, explaining and describing at B1.
\end{grammarlens}

\subsection*{Your turn}
\begin{itemize}
\raggedright
  \item \emph{Say it:} \emph{elhâm gereftan}
  \item \emph{Say it:} \emph{elhâm gereftan}, once more
  \item \emph{Say it:} a sentence of your own with \emph{elhâm gereftan}
  \item \emph{Recall:} say the Persian for to compose, then the Persian for to recite, then say \emph{elhâm gereftan} again
\end{itemize}

\begin{tcolorbox}[breakable,colback=teal!4,colframe=teal!35!black,title={Before you move on}]
What does \fa{الهام} \fa{گرفتن} mean? (\textquotedblleft{}To be inspired, to draw inspiration\textquotedblright{}.) Say it once more.
\end{tcolorbox}
\section[\texorpdfstring{\fa{تخیل} (takhayyol)}{تخیل (takhayyol)}]{\fa{تخیل} (takhayyol) — imagination}

\label{lesson:FA-C522-takhayyol}



\begin{quote}
Before the new one: say the Persian for to recite, then the Persian for to be inspired.
\end{quote}

\subsection*{You'll want to know: \fa{تخیل}}
\textbf{\fa{تخیل}} (\emph{takhayyol}) — \textquotedblleft{}imagination\textquotedblright{}.

\begin{quote}
\fa{بچه} \fa{ها} \fa{تخیل} \fa{قوی} \fa{دارند}.
\end{quote}

\begin{quote}
\emph{bachchehâ takhayyol-e qavi dârand.}
\end{quote}

\begin{quote}
Children have strong imaginations.
\end{quote}

The adjective is \emph{takhayyoli} 'imaginary'.

\begin{cousinweb}
\emph{takhayyol} is a loan, usually traced to Arabic \emph{taḫayyul}, 'imagining'. EtymDB has no entry for it.
\end{cousinweb}

\begin{grammarlens}[title={What you've built}]
One more word for telling, explaining and describing at B1.
\end{grammarlens}

\subsection*{Your turn}
\begin{itemize}
\raggedright
  \item \emph{Say it:} \emph{takhayyol}
  \item \emph{Say it:} \emph{takhayyol}, once more
  \item \emph{Say it:} a sentence of your own with \emph{takhayyol}
  \item \emph{Recall:} say the Persian for to recite, then the Persian for to be inspired, then say \emph{takhayyol} again
\end{itemize}

\begin{tcolorbox}[breakable,colback=teal!4,colframe=teal!35!black,title={Before you move on}]
What does \fa{تخیل} mean? (\textquotedblleft{}Imagination\textquotedblright{}.) Say it once more.
\end{tcolorbox}
